
\section*{4. ``Frobeniuseries''}
\phantomsection
\label{sga3:VIIA:4}
\addcontentsline{toc}{section}{4. ``Frobeniuseries''}

Fix a prime number \(p\), and let
\((\mathbf{Sch}/\mathbb F_p)\) be the category of schemes of
characteristic \(p\), that is, schemes over the prime field
\(\mathbb F_p\). Following the general conventions of this seminar,
we identify \((\mathbf{Sch}/\mathbb F_p)\) with a subcategory of
\((\mathbf{Sch}/\mathbb F_p)^\wedge\) by means of the functor \(h\)
of \SGARefLink{sga3:i:subsection:1-1}{Expos\'e I, 1.1}. We likewise use the isomorphism
\[
  \operatorname{Hom}(h_X,F)\xrightarrow{\sim}F(X)
\]
defined there to identify these sets whenever \(X\) is an
\(\mathbb F_p\)-scheme and \(F\) is an object of
\((\mathbf{Sch}/\mathbb F_p)^\wedge\).

\medskip
\noindent\textbf{Notation 4.0.}
\phantomsection
\label{sga3:VIIA:4.0}
\textit{If \(T\) is an \(\mathbb F_p\)-scheme, a \(T\)-functor is a
morphism \(q:F\to T\) in
\((\mathbf{Sch}/\mathbb F_p)^\wedge\) with target \(T\). For every
\(T\)-scheme \(r:X\to T\), the set of \(T\)-morphisms \(X\to F\),
that is, of \(\mathbb F_p\)-morphisms \(s:X\to F\) satisfying
\(q\circ s=r\), is denoted}
\[
  q(r),\qquad q(X/T),\qquad F(r),\qquad\text{or}\qquad F(X/T).
\]
\textit{It may also be denoted \(F(X)\) when no confusion with
\(\operatorname{Hom}(h_X,F)\) can result.}

\medskip
\noindent\textbf{4.1.}
\phantomsection
\label{sga3:VIIA:4.1}
For every scheme \(S\) of characteristic \(p\), let
\(\operatorname{fr}(S)\), or simply \(\operatorname{fr}\), denote the
endomorphism of \(S\) which is the identity on the underlying
topological space and which sends a section \(x\) of \(\mathcal O_S\)
over an open subset to \(x^p\).

The assignment
\[
  \operatorname{fr}:S\longmapsto\operatorname{fr}(S)
\]
is an endomorphism of the identity functor of
\((\mathbf{Sch}/\mathbb F_p)\).
This has the following consequences. If \(E\) is an
\(\mathbb F_p\)-functor, the assignment which sends every
\(\mathbb F_p\)-scheme \(S\) to the endomorphism
\(E(\operatorname{fr}(S))\) of \(E(S)\) is a functorial endomorphism
of \(E\), denoted \(\operatorname{fr}(E)\) or simply
\(\operatorname{fr}\). This notation is compatible with the
identification of \((\mathbf{Sch}/\mathbb F_p)\) with a subcategory
of \((\mathbf{Sch}/\mathbb F_p)^\wedge\). Moreover,
\(E\mapsto\operatorname{fr}(E)\) is itself an endomorphism of the
identity functor of \((\mathbf{Sch}/\mathbb F_p)^\wedge\), again
denoted \(\operatorname{fr}\).

For every \(\mathbb F_p\)-scheme \(S\) and every \(S\)-functor
\(q:X\to S\), let \(X^{(p/S)}\), or \(X^{(p)}\), denote the pullback
of \(X\) along \(\operatorname{fr}(S)\):
\begin{center}
  \phantomsection
  \label{sga3:VIIA:4.1-frobenius-basechange-square}
  \begin{tikzcd}[column sep=large,row sep=large]
X^{(p/S)}
  \arrow[r,"\operatorname{pr}_X"]
  \arrow[d]
&
X
  \arrow[d,"q"]
\\
S
  \arrow[r,"\operatorname{fr}(S)"]
&
S\,.
\end{tikzcd}

\end{center}

The commutative square
\begin{center}
  \phantomsection
  \label{sga3:VIIA:4.1-absolute-frobenius-square}
  \begin{tikzcd}[column sep=large,row sep=large]
X
  \arrow[r,"\operatorname{fr}(X)"]
  \arrow[d,"q"']
&
X
  \arrow[d,"q"]
\\
S
  \arrow[r,"\operatorname{fr}(S)"]
&
S
\end{tikzcd}

\end{center}
induces an \(S\)-morphism
\[
  \operatorname{Fr}(X/S):X\longrightarrow X^{(p/S)}
\]
such that
\[
  \operatorname{fr}(X)
  =
  \operatorname{pr}_X\circ\operatorname{Fr}(X/S):
\]
\begin{center}
  \phantomsection
  \label{sga3:VIIA:4.1-relative-frobenius-factorization}
  \begin{tikzcd}[column sep=large,row sep=large]
X
  \arrow[drr,bend left=24,"\operatorname{fr}(X)"]
  \arrow[dr,"\operatorname{Fr}(X/S)"]
  \arrow[ddr,"q"']
\\
&
X^{(p/S)}
  \arrow[r,"\operatorname{pr}_X"]
  \arrow[d]
&
X
  \arrow[d,"q"]
\\
&
S
  \arrow[r,"\operatorname{fr}(S)"]
&
S\,.
\end{tikzcd}

\end{center}
We call \(\operatorname{Fr}(X/S)\) the \emph{Frobenius morphism of
\(X\) relative to \(S\)}. Plainly the assignment
\[
  \operatorname{Fr}:X\longmapsto\operatorname{Fr}(X/S)
\]
is a functorial homomorphism.

\footnote{The editors expanded the original in the remainder of this
paragraph.}
Let \(r:T\to S\) be an \(S\)-scheme. For every
\[
  \varphi\in X(r)=\operatorname{Hom}_S(T,X),
\]
one has the commutative diagram
\begin{center}
  \phantomsection
  \label{sga3:VIIA:4.1-point-phi-diagram}
  {\small
\begin{tikzcd}[column sep=normal,row sep=large]
X
  \arrow[r,"\operatorname{Fr}(X/S)"]
  \arrow[dr,"q" description]
&
X^{(p/S)}
  \arrow[r,"\operatorname{pr}_X"]
  \arrow[d,"q^{(p/S)}"]
&
X
  \arrow[d,"q"]
\\
T
  \arrow[u,"\varphi"]
  \arrow[r,"r"]
&
S
  \arrow[r,"\operatorname{fr}(S)"]
&
S\,.
\end{tikzcd}}

\end{center}
By the definition of \(X^{(p/S)}\) as a fiber product,
\(\operatorname{pr}_X\) induces a bijection
\[
 X^{(p/S)}(r)
 =
 \operatorname{Hom}_S(T,X^{(p/S)})
 \xrightarrow{\sim}
 \operatorname{Hom}_S(T,X)
 =
 X(\operatorname{fr}(S)\circ r).
\]
On the other hand,
\[
  r\circ\operatorname{fr}(T)
  =
  \operatorname{fr}(S)\circ r,
\]
because \(\operatorname{fr}\) is an endomorphism of the identity
functor. Hence
\[
  \operatorname{Fr}(X/S)(r):
  X(r)\longrightarrow X^{(p/S)}(r)
\]
is characterized by commutativity of
\begin{center}
  \phantomsection
  \label{sga3:VIIA:4.1-pointwise-frobenius-square}
  \begin{tikzcd}[column sep=huge,row sep=large]
X(r)
  \arrow[r,"\operatorname{Fr}(X/S)(r)"]
  \arrow[d,"X(\operatorname{fr}(T))"']
&
X^{(p/S)}(r)
  \arrow[d,"\wr"]
\\
X(r\circ\operatorname{fr}(T))
  \arrow[r,equal]
&
X(\operatorname{fr}(S)\circ r)\,.
\end{tikzcd}

\end{center}

For example, if \(X\) is the subscheme of \(S\) defined by a
quasi-coherent ideal \(\mathcal I\), then \(X^{(p)}\) is the subscheme
defined by the ideal \(\mathcal I^{\{p\}}\) generated by the \(p\)-th
powers of local sections of \(\mathcal I\). Moreover,
\(\operatorname{Fr}(X/S)\) is the canonical immersion
\[
  \operatorname{Spec}(\mathcal O_S/\mathcal I)
  \longrightarrow
  \operatorname{Spec}(\mathcal O_S/\mathcal I^{\{p\}}).
\]
\footnote{The French source prints
\(\operatorname{Spec}(\mathcal O_X/\mathcal I)\) for the source of
this immersion. Since \(X\subset S\) is defined by \(\mathcal I\), the
intended source is
\(\operatorname{Spec}(\mathcal O_S/\mathcal I)=X\), as used here.}

\medskip
\noindent\textbf{4.1.1.}
\phantomsection
\label{sga3:VIIA:4.1.1}
\footnote{The editors expanded the original in this paragraph.}
Let \(t:T\to S\) be a base change and put
\[
  X_T=X\times_{q,t}T.
\]
Consider the pullback of \(X_T\) along \(\operatorname{fr}(T)\):
\begin{center}
  \phantomsection
  \label{sga3:VIIA:4.1.1-basechange-rectangle}
  \begin{tikzcd}[column sep=large,row sep=large]
(X_T)^{(p/T)}
  \arrow[r]
  \arrow[d]
&
X_T
  \arrow[r]
  \arrow[d]
&
X
  \arrow[d,"q"]
\\
T
  \arrow[r,"\operatorname{fr}(T)"]
&
T
  \arrow[r,"t"]
&
S\,.
\end{tikzcd}

\end{center}
Since
\[
  t\circ\operatorname{fr}(T)
  =
  \operatorname{fr}(S)\circ t,
\]
\((X_T)^{(p/T)}\) identifies with the pullback of \(X^{(p/S)}\)
along \(t\). Equivalently, there is a canonical isomorphism
\[
  X_T^{(p/T)}
  \xrightarrow{\sim}
  \bigl(X^{(p/S)}\bigr)_T.
\]
Under this identification,
\(\operatorname{Fr}(X_T/T)\) identifies with the pullback
\(\operatorname{Fr}(X/S)_T\).

\medskip
\noindent\textbf{4.1.1.1.}
\phantomsection
\label{sga3:VIIA:4.1.1.1}
In particular, if \(S=\operatorname{Spec}\mathbb F_p\), then
\(X^{(p/S)}=X\) and
\(\operatorname{Fr}(X/S)=\operatorname{fr}(X)\). Consequently,
\(X_T^{(p/T)}\) identifies with \(X_T\), and
\(\operatorname{Fr}(X_T/T)\) with
\(\operatorname{fr}(X)_T\).

For example, if \(E\) is a set and \(E_T\) the constant \(T\)-scheme
of type \(E\), then
\[
  E_T^{(p/T)}\simeq E_T,
  \qquad
  \operatorname{Fr}(E_T/T)\simeq\operatorname{id}_{E_T}.
\]

\medskip
\noindent\textbf{4.1.2.}
\phantomsection
\label{sga3:VIIA:4.1.2}
The functor \(X\mapsto X^{(p/S)}\) plainly commutes with products.
It therefore sends an \(S\)-group \(G\) to an \(S\)-group
\(G^{(p/S)}\). Since \(\operatorname{Fr}\) is functorial,
\[
  \operatorname{Fr}(G/S):G\longrightarrow G^{(p/S)}
\]
is a homomorphism of \(S\)-groups. We denote its kernel by
\(\operatorname{Fr}G\).

If \(r:T\to S\) is an \(S\)-scheme, the
\hyperref[sga3:VIIA:4.1-pointwise-frobenius-square]{pointwise
Frobenius diagram} shows that the value
of \(\operatorname{Fr}G\) at \(r\) is the kernel of
\[
  G(\operatorname{fr}(T)):
  G(r)\longrightarrow G(r\circ\operatorname{fr}(T)).
\]

When \(T=I_R\) is the scheme of dual numbers over an \(S\)-scheme
\(R\), the absolute Frobenius of \(I_R\) factors as
\begin{center}
  \phantomsection
  \label{sga3:VIIA:4.1.2-dual-numbers-factorization}
  \begin{tikzcd}[column sep=large]
I_R
  \arrow[r,"\mathrm{can.}"]
&
R
  \arrow[r,"\operatorname{fr}(R)"]
&
R
  \arrow[r,"s"]
&
I_R,
\end{tikzcd}

\end{center}
where \(s\) is the zero section. It follows that
\((\operatorname{Fr}G)(I_R)\) contains the kernel
\(\operatorname{Lie}(G/S)(R)\) of
\[
  G(s):G(I_R)\longrightarrow G(R),
\]
and hence
\[
  \operatorname{Lie}(G/S)
  =
  \operatorname{Lie}(\operatorname{Fr}G/S).
\]

